\begin{table}[H]\centering\renewcommand{\arraystretch}{1.2}
\begin{tabular}{|c|c|c|}\hline Номер запуску & Одне ядро, с & Чотири ядра, с\\\hline
1 & $2{,}285847$ & $0{,}602125$\\\hline
2 & $2{,}305389$ & $0{,}646983$\\\hline
3 & $2{,}308227$ & $0{,}624484$\\\hline
Середнє & $2{,}299821$ & $0{,}624531$\\\hline\end{tabular}
\caption{Час виконання для $N=1000$}
\end{table}
Для $N=1000$ коефіцієнт прискорення та ефективність становлять
\begin{equation} K_p=\frac{t_1}{t_4}=\frac{2{,}299821}{0{,}624531}=3{,}682,\qquad \eta=\frac{K_p}{4}\cdot100\%\approx 92{,}1\%.\end{equation}
\begin{table}[H]\centering\renewcommand{\arraystretch}{1.2}
\begin{tabular}{|c|c|c|}\hline Номер запуску & Одне ядро, с & Чотири ядра, с\\\hline
1 & $3{,}978902$ & $1{,}107869$\\\hline
2 & $3{,}987936$ & $1{,}158707$\\\hline
3 & $4{,}066953$ & $1{,}132493$\\\hline
Середнє & $4{,}011264$ & $1{,}133023$\\\hline\end{tabular}
\caption{Час виконання для $N=1200$}
\end{table}
Для $N=1200$ коефіцієнт прискорення та ефективність становлять
\begin{equation} K_p=\frac{t_1}{t_4}=\frac{4{,}011264}{1{,}133023}=3{,}540,\qquad \eta=\frac{K_p}{4}\cdot100\%\approx 88{,}5\%.\end{equation}
Величина $t_1$ є часом тієї самої чотирипотокової програми на одному
CPU, а не часом іншої послідовної реалізації. Результати залежать від
оптимізації компілятора, частоти, кешів і фонового навантаження.
Додаткові запуски на CPU 0--3 використовували лише два фізичні ядра з SMT;
їх не включено до розрахунку прискорення на чотирьох ядрах.
